\chapter{Overview: three geometries, one method}
\label{chap:overview}

% Architectural overview of the whole development. No formal statements are
% made here (no \lean pins); every result named below is stated and proved in
% its own chapter. Formalization map: Gluck/Basic.lean (umbrella import).

This blueprint documents a complete, axiom-clean formalization of the
\emph{converse to the four-vertex theorem} on the three two-dimensional space
forms: the Euclidean plane $\mathbb{E}^2$ ($K=0$), the round sphere
$\mathbb{S}^2$ ($K=+1$), and the hyperbolic plane $\mathbb{H}^2$ ($K=-1$).
In each geometry the converse takes the same shape: a prescribed curvature
profile $\kappa\colon\mathbb{R}\to\mathbb{R}$ satisfying a four-vertex-type
hypothesis is realized \emph{exactly} as the (geodesic) curvature of a simple
closed curve.

\section{The capstone theorems}

Each geometry carries a \emph{positive} converse (curvature bounded below by
the geometry's escape threshold) and a \emph{mixed-sign} (Dahlberg) converse
in which the minima are unconstrained:

\begin{center}
\begin{tabular}{lll}
 & positive curvature & mixed-sign curvature \\
\hline
$\mathbb{E}^2$ & \texttt{Gluck.gluck\_converse} & \texttt{Gluck.dahlberg\_converse} \\
$\mathbb{S}^2$ & \texttt{Gluck.spherical\_gluck\_converse} & \texttt{Gluck.spherical\_dahlberg\_converse} \\
$\mathbb{H}^2$ & \texttt{Gluck.hyperbolic\_gluck\_converse} & \texttt{Gluck.hyperbolic\_dahlberg\_converse} \\
\end{tabular}
\end{center}

All six conclusions are \emph{exact}: the witness curve realizes $\kappa$ on
the nose, with no reparametrization. In addition the $K$-generic
space-form layer states the unified capstones
\texttt{Gluck.SpaceForm.gluck\_converse} and
\texttt{Gluck.SpaceForm.dahlberg\_converse} (with the $\mathbb{H}^2$
instance \texttt{Gluck.SpaceForm.dahlberg\_converse\_hyperbolic}), and the
Euclidean development also proves the arc-length converse
\texttt{Gluck.arcLength\_converse}.

The geometry enters through the escape threshold: in $\mathbb{E}^2$ positivity
suffices; in $\mathbb{H}^2$ closed curves require geodesic curvature above the
horocycle value, so the positive converse assumes $\kappa>1$ and the mixed
converse requires the two maxima (only) to exceed an escape level $c>1$, while
the minima may be arbitrarily negative.

\section{Architecture}

Everything is modeled on plane curves: $\mathbb{S}^2$ and $\mathbb{H}^2$ are
realized on the open unit disk $\{\lVert z\rVert<1\}\subset\mathbb{C}$
(stereographic and Poincar\'e models), and the geometries differ only in the
\emph{realization predicate} coupling the tangent angle to the prescribed
curvature. Consequently the Euclidean chord, step-profile and winding
machinery is shared substrate for all three geometries --- embeddedness is a
property of the plane curve, independent of the ambient metric.

\begin{verbatim}
Common foundations (curves, curvature, winding)
        |
        v
Euclidean closing, simplicity, step machinery
        |
        +--------------> Sphere pipeline (flow engine)
        |                       |
        |                       v
        +--------------> SpaceForm reuse layer (epsilon-generic)
        |                       |
        |                       v
        +--------------> Hyperbolic pipeline (arc-length engine)
\end{verbatim}

The reuse is genuinely cross-cutting rather than tree-shaped, and this
blueprint states it as the implementation has it:
\begin{itemize}
\item the \emph{SpaceForm} layer (Part~4) is the spherical flow engine
  transported to an ambient-curvature parameter $K\in\{+1,-1\}$, re-deriving
  the positive spherical converse and producing the positive hyperbolic one;
\item the \emph{Hyperbolic} pipeline (Part~5) is a separate arc-length
  reconstruction engine built on the SpaceForm definitions --- $\mathbb{H}^2$
  has no metric rescaling, so the period of the reconstruction is
  co-constructed and the engine works in Euclidean arc length;
\item the mixed space-form capstone
  \texttt{Gluck.SpaceForm.dahlberg\_converse} \emph{assembles} the
  spherical and hyperbolic mixed capstones by reduction, rather than proving a
  third independent construction.
\end{itemize}

\section{Formalization map}

The source tree mirrors the part structure of this blueprint:
\texttt{Gluck/Curve.lean}, \texttt{Gluck/Curvature.lean},
\texttt{Gluck/Winding.lean} (Part~1);
\texttt{Gluck/Euclidean/} (Part~2);
\texttt{Gluck/Sphere/} (Part~3);
\texttt{Gluck/SpaceForm/} (Part~4);
\texttt{Gluck/Hyperbolic/} with its \texttt{ArcLength/}, \texttt{MixedSign/}
and \texttt{Family/} sub-developments (Part~5).
The umbrella module \texttt{Gluck/Basic.lean} imports all four geometries, so
\texttt{import Gluck} provides every capstone and \texttt{lake build} verifies
the entire development. All capstones depend only on the three standard axioms
(\texttt{propext}, \texttt{Classical.choice}, \texttt{Quot.sound}).
